\section{Discussion}
\label{sec:discussion}

Our findings establish directional decomposition not merely as an analytical probe, but as a lens that illuminates fundamental mechanisms of Transformer computation across their lifecycle.

\paragraph{Geometric Duality of Representation Evolution.}
Standard formulations treat Transformer layer depth as incremental vector accumulation $\mathbf{x} + \Delta$. Our decomposition reveals that this accumulation executes two fundamentally distinct geometric functions: \textit{magnitude modulation} ($\Delta_\parallel$) and \textit{semantic steering} ($\Delta_\perp$). Perpendicular updates rotate representations into new orthogonal subspaces where factual associations and syntactic distinctions reside; because these semantic coordinates are finely calibrated, $\Delta_\perp$ is hyper-fragile to perturbation. Conversely, parallel updates act as adaptive gain controllers, amplifying existing feature trajectories without altering their directional meaning. The extensive tolerance basin of $\Delta_\parallel$ proves that Transformers possess an inherent resilience to scale variations along their representational trajectories, explaining why coarse semantic intent persists even under substantial magnitude disruption.

\paragraph{Rethinking Model Compression: Beyond Isotropic $L_2$ Error.}
A ubiquitous convention in post-training quantization and pruning is the minimization of isotropic reconstruction error, $\min \|\Delta - \widehat{\Delta}\|_2^2$. However, isotropic $L_2$ error implicitly assumes uniform sensitivity across all dimensions of $\mathbb{R}^d$. Our directional framework exposes the flaw in this assumption: approximation errors aligned with the update ($\Delta_\parallel^{\mathrm{err}}$) fall within the model's natural tolerance basin, whereas orthogonal deviations ($\Delta_\perp^{\mathrm{err}}$) directly distort semantic trajectories. This resolves the long-standing mystery of why compression methods exhibiting comparable global $L_2$ error often diverge dramatically in downstream performance. Measuring compression distortion through directional decomposition establishes $\Delta_\perp^{\mathrm{err}}$ as a causal, high-fidelity predictor of downstream degradation, suggesting that future compression objectives should penalize directional rotation rather than uniform distance.

\paragraph{Architectural Implications and Pretraining Dynamics.}
Why do unconstrained Transformers consistently allocate substantial parameter capacity to redundant parallel scaling within attention? In standard self-attention, token mixing is structurally entangled with self-representation amplification. Our from-scratch pretraining results demonstrate that suppressing parallel updates relieves attention from redundant scaling duties, guiding optimization to dedicate its capacity toward directional contextual steering. This architectural inductive bias accelerates convergence, lowering validation-loss trajectories and improving downstream generalization across model scales from 300M to 2.7B parameters. These findings provide principled geometric support for emerging architectures that decouple contextual routing from magnitude modulation (such as gated attention and query-key normalization), indicating that directional orthogonality is a promising inductive bias for next-generation foundation models.
